\chapter{Тактирование, временной анализ и разные тактовые домены}
\label{ch:clocks}

Синхронная схема работает правильно, только если каждый сигнал успевает дойти
от одного триггера до другого за время одного такта. В главе~\ref{ch:sync} мы
вывели условие для максимальной частоты. Теперь посмотрим, как это проверяет
программа и что делать, если у схемы больше одного тактового сигнала или
сигнал приходит снаружи.

\section{Условие времени в ПЛИС}

Для пары триггеров, между которыми стоит комбинационная логика, должно
выполняться условие

\[
T \;\ge\; t_{cq} + t_{\text{лог}} + t_{\text{пров}} + t_{su},
\]

где $t_{cq}$ --- задержка выхода первого триггера после фронта, $t_{\text{лог}}$ ---
задержка логики (LUT), $t_{\text{пров}}$ --- задержка соединений и $t_{su}$ --- время
установки второго триггера. В ПЛИС задержка проводов составляет значительную
часть общего времени, поэтому её невозможно оценить, пока схема не размещена
на кристалле.

Именно поэтому после размещения и трассировки nextpnr выполняет
\term{статический временной анализ}: он перебирает все пути между
триггерами, складывает задержки и сравнивает с периодом такта. Разница
между периодом и длиной самого длинного пути --- \term{запас} (slack). Он
положителен, если схема укладывается, и отрицателен, если нет. Самый длинный
путь называют \term{критическим}, а предельную частоту, которую даёт его
задержка, --- $f_{max}$.

\begin{figure}[H]
\centering
\begin{tikzpicture}[font=\sffamily\small, x=1cm, y=1cm]
  \draw[->] (-0.2,0) -- (11,0) node[right] {время};
  \foreach \x/\t in {0/{фронт 1}, 8/{фронт 2}} {
    \draw[FpgaBlue, thick] (\x,-0.3) -- (\x,2.8); \node at (\x,3.1) {\t};
  }
  \fill[FpgaLight, draw=FpgaBlue] (0,0.2) rectangle (1.2,0.8) node[midway] {$t_{cq}$};
  \fill[FpgaLight!70!white, draw=FpgaBlue] (1.2,0.2) rectangle (3.4,0.8) node[midway] {$t_{\text{лог}}$};
  \fill[orange!20, draw=FpgaOrange] (3.4,0.2) rectangle (6.6,0.8) node[midway] {$t_{\text{пров}}$};
  \fill[FpgaLight, draw=FpgaBlue] (6.6,0.2) rectangle (7.5,0.8) node[midway] {$t_{su}$};
  \fill[green!20, draw=FpgaGreen] (7.5,0.2) rectangle (8,0.8);
  \draw[<->] (0,1.6) -- node[above] {период $T$} (8,1.6);
  \node[text=FpgaGreen!50!black] at (7.75,-0.6) {запас};
  \draw[->, FpgaGreen!50!black] (7.75,-0.4) -- (7.75,0.15);
\end{tikzpicture}
\caption{Путь между двумя триггерами. Если сумма задержек укладывается в период, остаётся положительный запас}
\end{figure}

\subsection{Как читать отчёт}

Откройте вкладку «Кристалл» и нажмите «Разместить и собрать». После
сборки в правой колонке появится результат: максимальная достижимая
частота для каждого тактового сигнала и сравнение с требуемой. На рис.~\ref{fig:n-chip}
сигнал \code{clk} (27~МГц) имеет $f_{max}$ около 384~МГц, то есть запас примерно
в 14 раз.

\shot{app-chip}{Отчёт nextpnr в Yosis: достижимая частота 384,47~МГц при требуемых 27~МГц}{fig:n-chip}

Если $f_{max}$ меньше нужной частоты, схема на плате будет работать
\emph{неправильно, но без сообщений об ошибке}: триггеры будут
захватывать значения, которые ещё не успели установиться. Поэтому после
каждой сборки нужно смотреть на отчёт. Лечить нехватку времени можно так:
разбить длинную логику конвейером (глава~\ref{ch:sync}), упростить выражение,
уменьшить число входов, разрешить синтезатору дублировать ячейки.

\begin{infobox}[Hold-проверка]
Кроме запаса по установке (setup) бывает запас по удержанию (hold): новое
значение не должно прийти \emph{слишком рано} после фронта, пока триггер ещё
удерживает старое. В ПЛИС этот запас обеспечивает сама трассировка и
общая тактовая сеть, поэтому обычно он заботит инженеров не так сильно.
\end{infobox}

\section{Не делите частоту --- используйте разрешение}

Начинающие часто «делят частоту» так: триггер переключается каждый
такт и его выход используют как новый тактовый сигнал для другой части схемы.
Это \emph{плохая} практика. Такой «такт» идёт не по специальной тактовой
сети, а по обычным проводам, из-за чего он приходит к триггерам с
разной задержкой, а анализ времени теряет надёжность.

Правильный способ --- оставить один тактовый сигнал на всю схему, а медленные
события отмечать \emph{импульсом разрешения} (clock enable). Делитель
считает такты и раз в $N$ тактов выдаёт короткий импульс \code{tick}, и только
тогда триггеры других частей схемы меняют значение.

\begin{vcode}
reg [24:0] cnt = 0;
reg        tick = 0;                      // один такт раз в секунду
always @(posedge clk) begin
    if (cnt == 27_000_000 - 1) begin
        cnt  <= 0;
        tick <= 1;
    end else begin
        cnt  <= cnt + 1;
        tick <= 0;
    end
end

always @(posedge clk)                      // тот же clk!
    if (tick) led <= ~led;                 // меняется раз в секунду
\end{vcode}

Этот приём вы использовали в проектах главы~\ref{ch:examples}. Сигнал
\code{tick} длится ровно один такт, но управляет любым числом триггеров, а все
они остаются синхронными с одним \code{clk}.

\section{PLL: умножение частоты}

На плате Tang Nano 20K стоит кварцевый генератор на 27~МГц, но
это не предел. Внутри кристалла есть \term{PLL} (phase-locked loop,
фазовая автоподстройка частоты): блок, который из входной частоты
$f_{\text{вх}}$ делает другую по формуле

\[
f_{\text{вых}} = f_{\text{вх}}\cdot\frac{M}{N},
\]

где $M$ и $N$ --- целые множители умножения и деления (точные
допустимые диапазоны есть в документации на кристалл). Так из 27~МГц можно
получить, например, 108~МГц для быстрой логики или 74{,}25~МГц для
видеосигнала HDMI. PLL также очищает тактовый сигнал от дрожания и может
вырабатывать несколько фаз.

PLL подключают как готовый примитив (в Gowin он называется \code{rPLL}), а
его выход идёт на глобальную тактовую сеть. После этого схема за PLL
считается новым \term{тактовым доменом} со своей частотой.

\section{Несколько тактовых доменов}

Когда в схеме два разных тактовых сигнала, возникает новая проблема:
сигнал, пришедший из одного домена в другой, может появиться как раз
\emph{во время} фронта принимающего такта. Тогда нарушается условие
установки/удержания, и триггер уходит в метастабильное состояние
(глава~\ref{ch:ff}): его выход зависает между 0 и 1 и через неопределённое время
решает сам. Если неудачное значение попадёт сразу в логику, схема
сломается.

\subsection{Синхронизатор из двух триггеров}

Простое решение --- пропустить сигнал через цепочку из \emph{двух} триггеров
принимающего домена. Первый может стать метастабильным, но у него есть
целый такт на то, чтобы «решить», и второй получает уже устойчивое
значение. Вероятность того, что метастабильность переживёт целый такт,
пренебрежимо мала.

\vfile{\codepath/10_clocks/sync2ff.v}{code/10\_clocks/sync2ff.v}

\shot{n-sync-rtl}{Синхронизатор в Yosis: два триггера один за другим, дальше фильтр дребезга}{fig:n-sync-rtl}

Такой же синхронизатор нужен для любого асинхронного сигнала: кнопки,
сигнала от другой микросхемы (например, входа приёмника последовательного
интерфейса, глава~\ref{ch:uart}).

\begin{warnbox}
Синхронизатор из двух триггеров годится только для \emph{одного бита}. Если
передавать многоразрядное число (например, счётчик), биты могут
прийти с разной задержкой, и принимающая сторона на мгновение увидит
значение, которого не было. Для шин применяют другие приёмы: передают по
одному биту с квитированием (запрос---подтверждение), используют \emph{код
Грея}, где при любом изменении меняется только один бит
(глава~\ref{ch:codes}), или двухтактные FIFO, у которых указатели
пересылаются в коде Грея.
\end{warnbox}

\subsection{Дребезг контактов}

Механическая кнопка при нажатии не замыкается один раз: контакты
подпрыгивают несколько миллисекунд, и сигнал дребезжит между 0 и 1. Цифровая
схема воспринимает это как множество нажатий. Исправляет это счётчик, который
принимает новое значение, только если оно держится $N$ тактов подряд.

\shot{n-debounce-wave}{Фильтр дребезга: вход \code{in} прыгает 1-0-1-0-1, но выход \code{out} переключается только после того, как значение продержалось несколько тактов}{fig:n-debounce}

На диаграмме (рис.~\ref{fig:n-debounce}) видно две важные вещи. Сигнал \code{s} после
синхронизатора повторяет \code{in} с задержкой в два такта. А выход
\code{out} ждёт, пока \code{s} установится: ложные короткие импульсы он
игнорирует.

\section{Сброс}

Сигнал сброса тоже должен быть синхронным. Асинхронный сброс срабатывает
мгновенно, но его \emph{снятие} в случайный момент может нарушить время
установки у части триггеров и привести схему в разные состояния. Правильно
подавать сброс на триггеры синхронно, а внешнюю кнопку сброса пропускать
через тот же двухтактный синхронизатор. В примерах книги везде используется
синхронный сброс: он проще и надёжнее, а кристаллу такого типа не нужны
отдельные цепи.

\section{Вопросы и задачи}
\label{sec:tasks-clocks}

Ответы собраны в приложении~\ref{app:answers}.

\begin{taskbox}[Задачи к главе]
\begin{enumerate}
  \item Путь между триггерами: $t_{cq}=0{,}5$~нс, логика $4$~нс, провода $3$~нс,
        время установки $0{,}5$~нс. Какой должен быть период? Какова $f_{max}$?
  \item Тактовый период 10~нс, суммарная задержка пути 8,2~нс. Каков запас?
        Каким он станет при 150~МГц?
  \item Опишите делитель, который выдаёт импульс \code{tick} с частотой 1~кГц
        из 27~МГц. До какого значения должен считать счётчик?
  \item PLL получает 27~МГц, $M=4$, $N=1$. Какая частота на выходе?
  \item Почему нельзя передавать 8-разрядный счётчик из одного домена в другой
        через восемь отдельных синхронизаторов?
  \item Фильтру дребезга требуется отсеять помехи до 5~мс при тактовой частоте
        27~МГц. Чему равно число $N$?
\end{enumerate}
\end{taskbox}
